\section{Metode}

\subsection{Formulasi Model A (LP 2-Sektor, Gaya Shen)}

Model A merepresentasikan sel sebagai pengoptimal alokasi proteom dua sektor.
Variabel keputusan: fraksi proteom glikolitik $\phi_G$, fraksi respiratorik
$\phi_R$, dan fluks serapan glukosa per sektor $J_{\text{glc},i}$.

Fungsi tujuan:
\begin{equation}
    \max \; J_\text{ATP} = \gamma_G \cdot J_{\text{glc},G} + \gamma_R \cdot J_{\text{glc},R}
    \label{eq:obj}
\end{equation}

dengan kendala:
\begin{align}
    J_{\text{glc},i} &\leq V_i \cdot \phi_i, \quad i \in \{G, R\} \label{eq:cap} \\
    \phi_G + \phi_R &\leq \Phi \label{eq:prot} \\
    J_{\text{glc},G} + J_{\text{glc},R} &\leq g_\text{avail} \label{eq:glc} \\
    \phi_G, \phi_R, J_{\text{glc},i} &\geq 0 \label{eq:nn}
\end{align}

Persamaan (\ref{eq:cap}) membatasi fluks glukosa oleh kapasitas enzim,
(\ref{eq:prot}) adalah anggaran proteom, dan (\ref{eq:glc}) adalah
ketersediaan glukosa --- knob utama (bukan O$_2$, karena efek Warburg
terjadi saat O$_2$ tersedia). Ini merupakan \textit{linear program} (LP)
yang diselesaikan menggunakan \texttt{scipy.optimize.linprog} dengan
backend HiGHS.

\subsection{Formulasi Model B (5-Parameter Kukurugya)}

Model B memaksimalkan $J_\text{ATP}$ yang sama terhadap kendala
(\ref{eq:prot})--(\ref{eq:nn}), namun cukup kecil untuk memiliki
\textbf{solusi analitik titik-sudut}. Tiga kandidat optimum:

\begin{itemize}
    \item \textbf{Respirasi murni}: $\phi_R = \min(\Phi, g/V_R)$, $\phi_G=0$
    \item \textbf{Glikolisis murni}: $\phi_G = \min(\Phi, g/V_G)$, $\phi_R=0$
    \item \textbf{Campuran}: kedua kendala mengikat $\Rightarrow$ sistem $2\times 2$
\end{itemize}

Lima parameter ($\gamma_R$, $\gamma_G$, $V_R$, $V_G$, $\Phi$) diambil
dari Tabel~1 dan SI Kukurugya \cite{kukurugya2024} untuk tiga organisme:
\textit{E.\ coli}, \textit{S.\ cerevisiae}, dan sel mamalia.

\subsection{Antarmuka Bersama}

Kedua model menerima \texttt{ModelParams} yang sama dan mengembalikan
\texttt{ModelResult} yang sama (fenotipe, $J_\text{ATP}$, fraksi
fermentatif, rasio efisiensi). Ini memastikan perbandingan pada
pijakan yang setara.

\subsection{E1: Overlay Sumbu-Bersama}

Kedua model dijalankan pada grid glukosa identik ($g_\text{avail}$ dari
$10^{-6}$ hingga $2{,}5 \cdot V_G \cdot \Phi$, 300 titik) per organisme.
Keluaran diplot pada sumbu bersama: sumbu-x = $g_\text{avail}$, sumbu-y =
fraksi fluks fermentatif.

\subsection{E2: Kapasitas versus Realisasi}

Eksperimen inti. Definisikan fraksi pemanfaatan kapasitas $u_i \in (0, 1]$
per sektor (1 = operasi di V\textsubscript{max}; $<1$ = ada enzim idle).
Dua efisiensi dihitung pada parameter yang sama:

\begin{align}
    \text{Kapasitas:} &\quad (V \cdot \gamma)_i \quad \text{(gaya Kukurugya, } u=1\text{)} \label{eq:cap_eff} \\
    \text{Terealisasi:} &\quad u_i \cdot (V \cdot \gamma)_i \quad \text{(gaya Shen)} \label{eq:real_eff}
\end{align}

Margin efisiensi didefinisikan sebagai:
\begin{equation}
    m = u_R \cdot (V \cdot \gamma)_R - u_G \cdot (V \cdot \gamma)_G
    \label{eq:margin}
\end{equation}

Verdict membalik saat $m = 0$, yaitu pada $u_G = \rho \equiv (V\!\cdot\!\gamma)_R / (V\!\cdot\!\gamma)_G$.
Kami menyapu $u_G$ dari 1{,}0 ke 0{,}01 (500 titik) dengan $u_R = 1$.

\textbf{Kuantifikasi ketidakpastian (E2c):} Bootstrap 200 iterasi,
menyampel parameter dari rentang $\pm 30\%$ nilai sentral, menghitung
distribusi titik-balik, dan melaporkan 90\% CI.

\textbf{Dekomposisi atribusi (E2d):} Variasi $\pm 30\%$ pada $V_G$ dan
$V_R$ (menyimulasikan ``enzim mana yang dihitung'') untuk mengkuantifikasi
kontribusi relatif definisi enzim terhadap efek $u_G$.

\subsection{E3: Identifiabilitas (Sobol + Morris)}

Analisis sensitivitas global dengan keluaran \textbf{kontinu}: margin $m$
dari (\ref{eq:margin}) dan ambang glukosa kritis $g^* = V_R \cdot \Phi$.
Distribusi input: 6 parameter ($\gamma_R$, $\gamma_G$, $V_R$, $V_G$,
$\Phi$, $u_G$) disampling dari rentang ketidakpastian SI.

\begin{itemize}
    \item \textbf{Sobol} \cite{salib}: $N=1024$ basis, indeks $S_1$ dan $S_T$ (total-order)
    \item \textbf{Morris}: $N=512$ trajectories, $\mu^*$ dan $\sigma$
\end{itemize}

Parameter teratas menurut $S_T$ = pengukuran yang, bila dipastikan,
paling mengurangi ketidakpastian verdict = eksperimen penuntas.

\begin{figure}[t]
    \centering
    \includegraphics[width=\columnwidth]{F1_overlay.png}
    \caption{Bagan alir arsitektur: M2 (parameter) $\to$ M1 (solver) $\to$ M3 (E1+E2) $\to$ M4 (E3) $\to$ M5 (visualisasi). Pipeline satu-perintah: \texttt{python main.py}.}
    \label{fig:architecture}
\end{figure}
